\answersfor{ch:recon}

\begin{solution}{ex:recon:1}
Reproducing the visible input has a perfect solution, the identity $\mathrm{id}(\vect{x})=\vect{x}$,
whose loss is exactly zero; training heads towards it, and the identity returns every pixel unchanged,
so the residual $\vect{x}-\xhat$ is zero everywhere, anomalies included, and the score map is empty.
The failure is dangerous because it is silent: nothing crashes and nothing warns, and an empty score map
looks exactly like a scene with nothing unusual in it.
\end{solution}

\begin{solution}{ex:recon:2}
Directly. First the reconstruction:
\begin{align*}
\hat{x} &= 0.25 \times 0.80 + 0.75 \times 0.20\\
&= 0.20 + 0.15\\
&= 0.35.
\end{align*}
Then the residual:
\[
r = x - \hat{x} = 0.80 - 0.35 = 0.45.
\]
With \cref{eq:recon:shrink}, $w = 0.25$ so $1-w = 0.75$, and the honest surprise is
$x - c = 0.80 - 0.20 = 0.60$:
\[
r = (1-w)(x-c) = 0.75 \times 0.60 = 0.45.
\]
The two agree. The model reports $0.45/0.60 = 0.75$, that is 75\% of the honest surprise; the
quarter it lost is exactly the self-weight.
\end{solution}

\begin{solution}{ex:recon:3}
Predictions with the leave-one-out neighbour average (ends copy their only neighbour):
\begin{align*}
\hat{x}_1 &= x_2 = 2.0,\\
\hat{x}_2 &= \tfrac12(2.0 + 2.0) = \tfrac12(4.0) = 2.0,\\
\hat{x}_3 &= \tfrac12(2.0 + 2.0) = 2.0,\\
\hat{x}_4 &= \tfrac12(x_3 + x_5) = \tfrac12(2.0 + 0.4) = \tfrac12(2.4) = 1.2,\\
\hat{x}_5 &= \tfrac12(x_4 + x_6) = \tfrac12(2.0 + 2.0) = 2.0,\\
\hat{x}_6 &= x_5 = 0.4.
\end{align*}
Residuals:
\begin{align*}
r_1 &= 2.0-2.0 = 0, & r_4 &= 2.0 - 1.2 = 0.8,\\
r_2 &= 2.0-2.0 = 0, & r_5 &= 0.4 - 2.0 = -1.6,\\
r_3 &= 2.0-2.0 = 0, & r_6 &= 2.0 - 0.4 = 1.6.
\end{align*}
Scores: $0,\ 0,\ 0,\ 0.8,\ 1.6,\ 1.6$. The anomaly (pixel~5, a dark pixel) scores $1.6$, but pixel~6
ties with it. Pixel~6 sits at the end of the signal and has only one neighbour, the anomaly itself, so
its whole prediction is the anomalous value and its halo is as large as the anomaly's own score.
The halo is pixels~4 ($0.8$) and~6 ($1.6$). The lesson: near the edge of the valid data a prediction
depends on fewer neighbours and can be badly contaminated, which is why the training code erodes the
validity mask near borders (\cref{ch:masking}).
\end{solution}

\begin{solution}{ex:recon:4}
L2 is minimised by the mean. Add the values one at a time:
\begin{align*}
0.9 + 1.0 &= 1.9,\\
1.9 + 1.1 &= 3.0,\\
3.0 + 6.0 &= 9.0,
\end{align*}
and divide by 4: $9.0/4 = 2.25$.

L1 is minimised by any median. Sorted, the values are $0.9,\ 1.0,\ 1.1,\ 6.0$; the two middle values
are $1.0$ and $1.1$, and the midpoint is
\[
\tfrac12(1.0+1.1) = \tfrac12(2.1) = 1.05.
\]
(Any guess between $1.0$ and $1.1$ has two values above and two below, so the slope of the L1 loss is
zero there.)

The mean is pulled towards the anomaly. The three ordinary values average
$(0.9+1.0+1.1)/3 = 3.0/3 = 1.0$; the L2 guess $2.25$ is $2.25-1.05 = 1.2$ higher than the L1 guess
and more than double the typical ordinary value. A model trained with L2 would predict about $2.25$
behind this context and so report a residual of only $6.0-2.25 = 3.75$ at the anomaly, against
$6.0-1.05 = 4.95$ under L1.
\end{solution}

\begin{solution}{ex:recon:5}
Row~3 of $\vect{m}^{(1)}$ is $(0,1,0,1)$. At $(3,0)$, $m^{(1)}=0$: hidden in pass~1, so
$p^{(1)} = 1$ and the value comes from pass~1 ($i+j=3$ is odd, matching the pattern). At $(3,3)$,
$m^{(1)}=1$ and $m^{(2)}=0$: hidden in pass~2, so the value comes from pass~2 ($i+j=6$ is even).

No mixing. At any pixel $p^{(1)} = 1-m^{(1)}$ and $p^{(2)} = 1-m^{(2)} = m^{(1)}$. So
$p^{(1)}+p^{(2)} = 1$ and $p^{(1)}p^{(2)} = (1-m^{(1)})m^{(1)} = 0$, because $m^{(1)}$ is $0$ or $1$
and one of the two factors is then $0$. Two numbers that are each $0$ or $1$, sum to $1$ and multiply
to $0$ must be one $1$ and one $0$. \Cref{eq:recon:combine} then reads either
$\hat{x}^{(1)}\cdot1 + \hat{x}^{(2)}\cdot0 = \hat{x}^{(1)}$ or $\hat{x}^{(1)}\cdot0 +
\hat{x}^{(2)}\cdot1 = \hat{x}^{(2)}$: one pass, whole, never a blend.
\end{solution}

\begin{solution}{ex:recon:6}
(1) Token deletion at the encoder input: without it the hidden tokens' content would reach the
encoder and could be copied into the output, a self-weight $w>0$. (2) Two complementary passes at
inference: with one pass, the pixels that were visible in it would be scored on reconstructions made
while they were visible, so their anomalies would be shrunk by $(1-w)$, possibly to zero. (3) The loss
mask restricting the grade to hidden positions: if visible positions were graded, training would reward
copying them and pull the network towards the identity, teaching it a skill that hurts detection.
\end{solution}
